% !TeX program = lualatex
% Compile twice: lualatex 230.tex
% All references and prose are embedded; no BibTeX database or external inputs.
\documentclass[11pt,a4paper]{article}
\usepackage[margin=24mm,headheight=14pt]{geometry}
\usepackage{fontspec}
\setmainfont{Latin Modern Roman}
\setsansfont{Latin Modern Sans}
\setmonofont{Latin Modern Mono}
\usepackage{amsmath,amssymb,mathtools}
\usepackage{unicode-math}
\setmathfont{Latin Modern Math}
\usepackage{microtype}
\usepackage{enumitem,xurl,needspace}
\usepackage[dvipsnames]{xcolor}
\usepackage{hyperref}
\hypersetup{unicode=true,colorlinks=true,linkcolor=MidnightBlue,urlcolor=MidnightBlue,
 pdftitle={Research Notebook: Section 230},pdfauthor={Research notebook},bookmarksnumbered=true}
\usepackage{bookmark}
\usepackage{tocloft}
\setlength{\cftsecnumwidth}{3em}
\cftsetpnumwidth{2.5em}
\cftsetrmarg{4em}
\usepackage{titlesec}
\titleformat{\section}{\Large\bfseries\raggedright}{\thesection}{0.7em}{}
\usepackage{fancyhdr}
\pagestyle{fancy}\fancyhf{}
\fancyhead[L]{\small\sffamily Open Problems | Research notebook}
\fancyhead[R]{\small 230 | 27 September 2026}
\fancyfoot[C]{\thepage}
\setlength{\parindent}{0pt}
\setlength{\parskip}{5pt plus 1pt}
\setlength{\emergencystretch}{3em}
\setcounter{tocdepth}{1}
\setcounter{secnumdepth}{1}
\setlist[itemize]{leftmargin=3em,itemsep=5pt,topsep=4pt,parsep=0pt}
\allowdisplaybreaks
\newcommand{\N}{\mathbb{N}}
\newcommand{\Z}{\mathbb{Z}}
\newcommand{\Q}{\mathbb{Q}}
\newcommand{\R}{\mathbb{R}}
\newcommand{\C}{\mathbb{C}}
\newcommand{\F}{\mathbb{F}}
\newcommand{\A}{\mathbb{A}}
\newcommand{\PP}{\mathbb P}
\newcommand{\E}{\mathbb{E}}
\newcommand{\eps}{\varepsilon}
\newcommand{\dd}{\,\mathrm d}
\newcommand{\sref}[2]{\hyperlink{source:#1:#2}{[S#2]}}
\newcommand{\eref}[2]{\hyperlink{extra:#1:#2}{[E#2]}}
% Turn the next switch off for a shorter reading copy. The quotations preserve
% every exactTarget field, including qualifications omitted from a short summary.
\newif\ifshowcatalogue
\showcataloguetrue
\newcommand{\cataloguescope}[1]{\ifshowcatalogue
 \paragraph{Exact scope in the supplied catalogue.}
 \begin{quote}\small #1\end{quote}\fi}

% Independently compilable extraction; original section text retained.
\numberwithin{equation}{section}
\begin{document}
\setcounter{section}{229}
% ================================================================
% problemId: problem.shafarevich-freeness-conjecture
% source releaseRank: 230; source releaseStatus: open
% exactTarget SHA-256: 2c2ea1514c0f6e8cd45f65614bb655accf3208789149ddba3cfa77a4626e3c95
\clearpage
\section{\texorpdfstring{Shafarevich Freeness Conjecture over Qab}{Shafarevich Freeness Conjecture over Qab}}\label{230}
\subsection{Definitions and mathematical statement}
Let $\Q^{\mathrm{ab}}$ be the union, in $\overline{\Q}$, of all finite abelian extensions of $\Q$. The target is
\[
 \operatorname{Gal}(\overline{\Q}/\Q^{\mathrm{ab}})\cong\widehat F_\omega
\]
as profinite groups, where $\widehat F_\omega$ denotes the free profinite group on a countably infinite basis converging to the identity. Its universal property extends maps from that basis into any profinite group, when the images converge to the identity, uniquely to continuous homomorphisms. This topological basis convention matters for infinite rank. The claim is not freeness as an abstract discrete group or merely freeness of every finite quotient.
\par\smallskip\textit{Formulation and concept references:} \sref{230}{1}.
\cataloguescope{Let Qab be the maximal abelian extension of Q in an algebraic closure Qbar (equivalently the maximal cyclotomic extension). The absolute Galois group Gal(Qbar/Qab) is a free profinite group of countably infinite rank, in the standard profinite convention of a free basis converging to the identity.}
\subsection{Short English statement}
After adjoining every abelian extension of the rationals, is the remaining absolute Galois group free profinite of countably infinite rank?
% BEGIN RESEARCH ATTEMPT 230
\subsection{Research attempt}

\paragraph{Current frontier and source access.}
The author-hosted file at \sref{230}{1} did not load in this session.
The displayed group-theoretic target is self-contained, and the same
Shafarevich formulation is explicitly given in the Introduction of
Harbater--Stevenson. Their Corollary~2.8 distinguishes projectivity from
proper solvability of finite split embedding problems; their geometric
and local-field results are not a theorem for $\Q^{\rm ab}$.
\eref{230}{1}. For global fields, the fact that the maximal abelian
extension has cohomological dimension one is classical; Gille explicitly
records this in the concluding questions of his survey.
\eref{230}{3}, Section~7. No verified 2024--2026 result obtained in this
search supplies the missing unrestricted proper-solvability assertion.

\paragraph{Attack route.}
One route is to kill cohomological obstructions using the large supply of
cyclotomic extensions. A second must ensure that a lifted homomorphism
has the full required image, rather than merely being a weak lift. The
third would build a free basis directly, but would still need compatible
solutions to infinitely many finite constraints. We pursue the first two
routes explicitly. The calculations show that the first is effective at
killing Brauer classes and leaves the second, nonabelian image problem
untouched. They also exhibit infinitely many independent quadratic
extensions, so the countable-rank issue is not silently replaced by
finitely generated profinite freeness.

\paragraph{Attempt.}
Put $K=\Q^{\rm ab}$ and $\Gamma=\operatorname{Gal}(\overline\Q/K)$.
All field embeddings below are taken in the fixed algebraic closure.

\textbf{1. Explicit cyclotomic splitting of Brauer classes.}
Let $F$ be any intermediate field $K\subset F\subset\overline\Q$,
including a finite extension of $K$, and let $A$ be a central simple
$F$-algebra. Its multiplication constants and a certificate for central
simplicity descend to a number field $L\subset F$, after enlarging $L$
if necessary. Write $A_L$ for the descended algebra. Only finitely many
local invariants $\operatorname{inv}_v(A_L)\in\Q/\Z$ are nonzero.
For each nonarchimedean such $v$ over a rational prime $p$, let $h_v$
be the order of the invariant and let $f_v$ be the residue degree of
$L_v/\Q_p$. Choose
\[
 N_v=2f_vh_v,\qquad m_v=p^{N_v}-1.
\]
The local cyclotomic extension $\Q_p(\mu_{m_v})/\Q_p$ is unramified
of degree $N_v$. Indeed, $m_v$ is prime to $p$, and the order of $p$
modulo $p^{N_v}-1$ is exactly $N_v$: no smaller positive exponent can
make the smaller positive integer $p^r-1$ divisible by $p^{N_v}-1$.
Its intersection with $L_v$ is the unramified extension of degree
$f_v$, since $f_v\mid N_v$. Therefore
\begin{equation}\label{230:localdegree}
 [L_v(\mu_{m_v}):L_v]=N_v/f_v=2h_v.
\end{equation}
Restriction of a local Brauer class multiplies its invariant by the local
extension degree, so \eqref{230:localdegree} kills that invariant.
The local invariant rule and the global injectivity
$\operatorname{Br}(L)\to\bigoplus_v\operatorname{Br}(L_v)$ are the
class-field-theoretic inputs here. \eref{230}{2}, Chapters~III--IV and
VIII, Theorem~4.2.

Now put $L'=L(i,\{\mu_{m_v}\}_v)\subset F$; the inclusion follows
because $K$ contains all roots of unity. Each completion over a formerly
bad place has degree divisible by the degree used in
\eqref{230:localdegree}. Real places become complex, and restriction
cannot create a nonzero invariant at a place where the original one was
zero. Every local invariant of $A_L\otimes_LL'$ is therefore zero.
The global injection gives that it is split, so $A$ is split as well.
We have proved
\begin{equation}\label{230:brauer}
 \operatorname{Br}(F)=0\quad\text{for every }K\subset F\subset\overline\Q.
\end{equation}
This is an explicit form of a known class-field-theoretic phenomenon,
not a new solution of Shafarevich's conjecture.

The usual Brauer criterion for cohomological dimension, applied to all
finite extensions of a characteristic-zero field, identifies the
vanishing in \eqref{230:brauer} with $\operatorname{cd}(\Gamma)\leq1$.
Equivalently $\Gamma$ is projective as a profinite group. We use these
standard cohomological criteria, not an assertion that a single vanishing
$\operatorname{Br}(K)=0$ alone always suffices. The projectivity/\allowbreak
cohomological-dimension equivalence and the global-field case are recalled
in \eref{230}{1}, Section~2, and \eref{230}{3}, Section~7.

\textbf{2. The rank is genuinely countably infinite.}
Since $\overline\Q$ is countable, $\Gamma$ has only countably many open
subgroups and is countably based. To see that its rank cannot be finite,
choose distinct rational primes $p_1,p_2,\ldots$ and take their positive
square roots in $K$. Their classes in $K^\times/K^{\times2}$ are
linearly independent over $\F_2$. For if a nonempty finite product were
a square, there would be $b\in K$ with
\[
 b^2=\prod_{i\in I}\sqrt{p_i}=\sqrt a,
 \qquad a=\prod_{i\in I}p_i>1\text{ square-free}.
\]
Then $b=\pm a^{1/4}$ under our chosen complex embedding. The polynomial
$X^4-a$ is Eisenstein at any prime dividing $a$. The field
$\Q(a^{1/4})$ has degree four and is real, but is not normal over $\Q$,
since it does not contain the conjugate $ia^{1/4}$. Every finite subfield
of the abelian Galois extension $K/\Q$ is itself Galois over $\Q$.
This contradiction proves independence.

Kummer theory now gives quotients $(\Z/2\Z)^r$ of $\Gamma$ for every
$r$. No finite set can topologically generate all those quotients.
Together with countable basing this gives countably infinite profinite
rank, with the convergent-generating-set convention used in the target
and in \eref{230}{1}, Section~2. Merely observing that $\Gamma$ is a
subgroup of the countably based absolute Galois group would not by itself
have ruled out finite rank.

\textbf{3. Proper split embedding problems isolate the missing step.}
An embedding problem is a pair of continuous epimorphisms
\[
 \alpha:\Gamma\twoheadrightarrow A,\qquad
 \beta:B\twoheadrightarrow A,
\]
where $A,B$ are finite. A weak solution is $\lambda:\Gamma\to B$ with
$\beta\lambda=\alpha$; a proper solution is such a $\lambda$ that is
surjective. A split problem additionally has a homomorphic section of
$\beta$. Projectivity supplies weak solutions, not proper ones.
The criterion in \eref{230}{1}, Corollary~2.8, shows that the exact
Shafarevich target would follow from
\begin{equation}\label{230:proper}
 \text{every finite split embedding problem for }\Gamma
           \text{ has a proper solution.}
\end{equation}

Here is the concrete pullback step behind the criterion, which checks the
image requirement. Start with any finite problem and a weak solution
$\lambda_0$, and let $C=\lambda_0(\Gamma)\subset B$.
The restriction $\beta|_C:C\twoheadrightarrow A$ is onto. Form
\[
 E=B\times_A C=\{(b,c):\beta(b)=\beta(c)\}.
\]
Projection $E\to C$ splits by $c\mapsto(c,c)$, so the problem with
$\Gamma\twoheadrightarrow C$ supplied by $\lambda_0$ is split.
If \eqref{230:proper} holds, it has a surjective solution
$\Gamma\twoheadrightarrow E$. Projection $E\to B$ is surjective,
because $C\to A$ is surjective. Composing gives a proper solution of
the original problem. The countable-rank form of Iwasawa's criterion
then gives $\Gamma\cong\widehat F_\omega$.
This standard reduction is included to show exactly where surjectivity
enters; it is not a new freeness criterion. \eref{230}{1},
Corollary~2.8 and its alternative proof.

\textbf{4. Trying to replace properness by cohomology fails.}
For a split problem $B=N\rtimes A$, the section gives the weak solution
$\gamma\mapsto(1,\alpha(\gamma))$ immediately. Its image can miss all
of the kernel $N$. Producing a suitable nonabelian crossed homomorphism
into $N$ with sufficiently large image is the extra task; making an
abelian degree-two obstruction vanish says nothing by itself about that
image. The elementary example $A=1,B=\Z/3\Z$ already distinguishes a
trivial weak solution from a proper solution.

There is also a structural counterexample to the implication
``projective and countably infinite rank implies free profinite.''
The free pro-$2$ group of countably infinite rank has $2$-cohomological
dimension one and $\ell$-cohomological dimension zero for $\ell\ne2$,
so it is projective as a profinite group. It has no quotient $\Z/3\Z$,
and hence is not free profinite of countably infinite rank. Here
``free pro-$2$'' is a different category; this group is not proposed as
$\Gamma$. This diagnostic uses the same projectivity criterion recalled
in \eref{230}{1}, Section~2, and shows why the missing condition cannot
be dropped on group-theoretic grounds.

\paragraph{Stress test.}
The local-degree construction kills classes only after a finite
cyclotomic extension chosen from their actual local data; it does not
claim that one fixed cyclotomic extension kills every algebra. It treats
all finite extensions of $K$ and the archimedean places, as required by
the cohomological argument. The independent Kummer classes are
$\sqrt{p_i}$, not the rational primes $p_i$, which are already squares
in $K$.

More seriously, realizing each finite group separately would not
necessarily solve the embedding problems with a \emph{prescribed}
quotient map $\alpha$. Taking inverse limits of unrelated finite
realizations would not produce a compatible free basis. Patching
results over local or geometric fields need their own hypotheses; the
existence of roots of unity does not turn $\Q^{\rm ab}$ into one of
those fields. No largeness or rational-point property has been inserted
without proof. The split-proper statement \eqref{230:proper} remains
unproved for this field in this attempt.

\paragraph{Outcome.}
\textbf{Outcome: Conditional reduction.}

The attempt makes the cyclotomic annihilation of Brauer classes explicit,
verifies infinite countable rank through independent quadratic
extensions, and isolates proper split embedding problems as a precise
sufficient condition for the stated profinite freeness. These steps
rederive known structural reductions and do not constitute a claimed
advance in the unrestricted embedding problem.

The unresolved issue is a construction of surjective lifts for every
finite split problem with its prescribed quotient. Weak lifts,
cohomological vanishing and countable rank have all been separated from
that issue. No free basis or counterexample to the target is produced.
% END RESEARCH ATTEMPT 230
\subsection{Sources}
\begin{itemize}\raggedright
\item[\textnormal{[S1]}]\hypertarget{source:230:1}{} David Harbater, Shafarevich Conjecture (author-hosted article).
\url{https://www2.math.upenn.edu/~harbater/sc1.pdf}.
{\small\textit{Catalogue formulation source.}}
% BEGIN ADDED REFERENCES 230
\item[\textnormal{[E1]}]\hypertarget{extra:230:1}{}
D. Harbater and K. F. Stevenson, \emph{Local Galois theory in dimension
two: Second edition}, arXiv:0907.3082v1, 17 July 2009, Introduction;
Section~2, definitions and Corollary~2.8 with its alternative proof.
\url{https://arxiv.org/html/0907.3082v1}.
Consulted 27 September 2026; this is a separately retrieved source,
not a claim that the catalogue's inaccessible author-hosted PDF was read.
\item[\textnormal{[E2]}]\hypertarget{extra:230:2}{}
J. S. Milne, \emph{Class Field Theory}, version 4.03, 6 August 2020,
Chapters~III--IV (local invariants) and VIII, Theorem~4.2 (global Brauer
exact sequence), with Chapter~VII, Section~7.
\url{https://www.jmilne.org/math/CourseNotes/CFT.pdf}.
Consulted 27 September 2026.
\item[\textnormal{[E3]}]\hypertarget{extra:230:3}{}
P. Gille, \emph{Serre's conjecture II: a survey}, in \emph{Quadratic
forms, linear algebraic groups, and cohomology}, Developments in
Mathematics 18, Springer (2010), 41--56; consulted author manuscript,
Section~7, p.~15, on cohomological dimension of the abelian closure of a
global field. \url{https://www.math.uni-bielefeld.de/lag/man/326.pdf}.
Consulted 27 September 2026; historical status lists are not treated as
current frontier assertions.
% END ADDED REFERENCES 230
\end{itemize}

\end{document}
